\section{Introduction}
Repeated use resolves different directions of a measurement at different
rates. For a projective qubit measurement, a change of basis can be
resolved at order $N^{-1}$ after $N$ calls. In a strictly noisy family
the corresponding order is $N^{-1/2}$. When the probabilities and the
basis both vary, the interaction of these scales determines the size
of a reusable description. This paper gives a finite-distance analysis
of the full body of ordered binary qubit measurements.

Write the positive effect as
\[
 E_{b,x}=\tfrac12\bigl((1+b)I+x\cdot\sigma\bigr),
 \qquad |b|+|x|\le1,
\]
and the negative effect as $I-E_{b,x}$. The outcome labels are fixed.
The parameter set is the set of effects themselves: when $x=0$ there
is no direction coordinate, but the scalar probability $(1+b)/2$
remains observable. Put
\[
 p=\tfrac12(1+b+|x|),\qquad q=\tfrac12(1+b-|x|).
\]
Thus $0\le q\le p\le1$ are the two eigenvalues of the positive effect.

For memoryless devices of the same interface, let
\[
 d_N(\Phi,\Psi)=\sup_{\mathsf T}
       \|\mathsf T[\Phi]-\mathsf T[\Psi]\|_1.
\]
The same causal tester is used under both hypotheses. It may retain
an entangled reference, use arbitrary finite quantum memory and
feedback, and stop at a public time bounded by $N$. The norm is
unhalved, so the maximum distance is two. Every invocation applies
the same measurement, and the retained output includes both classical
and quantum registers of the tester.

\subsection{The metric and the covering law}
There are two spectral scales and one angular scale. Define
\begin{align*}
 T_N(s,t)&=|s-t|\sqrt{\frac{N}
   {\max\{s(1-s),t(1-t)\}+N^{-1}}},\\
 K_N(p,q)&=(p-q)\sqrt{\frac{N}
   {p(1-p)+q(1-q)+N^{-1}}}.
\end{align*}
For a second effect with spectrum $p',q'$ and nonzero vector $x'$,
let $h\in[0,\pi]$ be the angle between $x$ and $x'$. Set
\[
 H_N=T_N(p,p')+T_N(q,q')
       +\min\{K_N(p,q),K_N(p',q')\}h,
\]
with the last term zero whenever one vector vanishes.
Theorem~\ref{thm:biasedmetric75} gives absolute $c,C>0$ such that
\[
 c\min\{1,H_N\}\le d_N(\mathcal M_{b,x},\mathcal M_{b',x'})
                  \le\min\{2,CH_N\}
\]
for every pair and every $N\ge1$. At one use the distance is exactly
$|b-b'|+|x-x'|$, including reference assistance.

Let $\mathcal C_N(\delta)$ be the least number of radius-$\delta$
balls covering the full effect body, with centres allowed to be
arbitrary legal memoryless devices of the same binary interface.
Let $\mathcal C_N^{\rm rat}(\delta)$ require rational effects as
centres. Theorem~\ref{thm:biasedcover75} proves that there are absolute
constants $c_0,C_0,\delta_0>0$ such that
\[
 c_0\frac{N^2\log(N+2)}{\delta^4}
 \le \mathcal C_N(\delta)
 \le \mathcal C_N^{\rm rat}(\delta)
 \le C_0\frac{N^2\log(N+2)}{\delta^4},
 \qquad N\ge1,\quad 0<\delta\le\delta_0.
\]
Consequently the optimal fixed-length description has length
\[
 2\log_2N+\log_2\log(N+2)+4\log_2(1/\delta)+O(1).
\]
The remainder is uniform in horizon and in this small-error range.
Both spectral coordinates, including bias, and the direction are
charged in one reusable index. Theorem~\ref{thm:biasedcodec75}
constructs an exact rational encoder and decoder attaining this order.

The logarithm is a property of the whole effect body. Near a projective
measurement put $\alpha=1-p$ and $\beta=q$. A region with
$\alpha,\beta$ of order $t$ has two spectral resolutions and two
angular resolutions of order $\delta\sqrt{t/N}$. Each separated
dyadic region between $t\asymp N^{-1}$ and a fixed positive depth
therefore contributes order $N^2\delta^{-4}$ cells. There are order
$\log(N+2)$ such regions. The constructive upper bound sums the same
scales in rational spectral coordinates. This argument covers the
scalar and rank-deficient strata without imposing an Ahlfors
hypothesis on the full four-dimensional body.

\subsection{Proof structure}
The spectral comparison reduces aligned measurements to a controlled
binary classical channel. Testing an extreme eigenvector prevents an
angular change from hiding either spectral endpoint. For the angular
upper bound, a common rotation of the discarded Kraus coordinates
makes the cross-overlap of two measurement dilations a scalar multiple
of the identity. This scalar overlap persists under every common
adaptive network and gives a finite-use fidelity bound. A randomized
bias-removing operation and a tangential product probe supply matching
lower witnesses. The cutoff $N^{-1}$ follows from these finite-use
estimates, including zero eigenvalues.

These arguments distinguish a projective measurement, where $p=1$ and
$q=0$, from a one-sided rank-deficient measurement with $q=0<p<1$.
The former has angular scale $N$; the latter has scale
$\sqrt{Np/(1-p)}$ at a fixed $p<1$. Rank deficiency alone therefore
does not determine the rate. The lower witnesses may depend on the
pair of hypotheses. They establish metric separation and covering
lower bounds; a single estimator for the entire body is a separate
question.

The exact code uses rational spectral charts and signed-axis
stereographic charts. A rational Cartesian target can have an
irrational norm. The encoder selects its spectral layers and direction
digits by sign-sensitive rational comparisons, without computing
that norm. Each valid index decodes to a positive effect bounded by
the identity. The encoder is a finite algorithm with numerical grid
enumeration; its payload bound is distinct from the time and mutable
memory used to construct the codebook.

\subsection{Relation to the earlier results}
The full-body theorem contains the unbiased slice $b=0$, whose more
precise Euclidean geometry is retained in
Theorem~\ref{thm:ballmetric74}. On a compact identifiable
$k$-Ahlfors-regular subset $X$ of that slice,
Theorem~\ref{thm:weightedcover74} gives the small-error law
\[
 \mathcal C_N(X,\delta)\asymp_X
 \delta^{-k}\int_X
       \left[\frac{N}{1-|x|^2+N^{-1}}\right]^{k/2}\,d\mu(x).
\]
Its error cap depends on the stated regularity data. A two-dimensional
disk has depth exponent $\alpha=1=k/2$, so its integral is of order
$N\log(N+2)$; a three-dimensional ball has $\alpha=1<k/2$, and its
boundary layer gives order $N^2$. The full biased body's new logarithm
arises from a two-endpoint spectral corner. These two critical
phenomena concern different target families.

We also retain the complete proofs for strictly positive instrument
descriptions, singular preparations, coherent tensor families,
observable varying readouts and uniformly seizable families. Their
dimensions and error ranges are stated in their respective theorems.
The structural companion develops stochastic realization and finite
physical actions with fresh classical nondisturbing probes. Its
results remain part of the complete research edition.

Single-use and multiple-use measurement discrimination supply direct
antecedents. Fiur\'a\v sek--Mi\v cuda \cite{FM75} compare two-use
projective-qubit strategies, including adaptivity, entangled probes
and feed-forward. Sedl\'ak--Ziman \cite{SZ74} treat noisy single-use
measurements with unequal visibilities; Pucha\l a et al.
\cite{PPKK74,PPKKO72} give projective single-use and multiple-use
distance formulas. The finite-distance proofs here establish the
joint biased-body comparison and its covering consequences. The
comparison section records the precise roles of those antecedents.

Throughout, the input to the encoder is a supplied mathematical
description. Dimension, horizon and requested accuracy are public;
the target effect is charged. The decoded object is a quantum
measurement. Theorems about its description length concern neither
unknown-device sample complexity nor the hardware needed to implement
the operation on a physical quantum input.
